%%% FIELD corrected-exact-projection-definition
For a multinomial count vector \(u\) of total \(N\geq 1\), let \(\ell_p(u)=\sum_j u_j\log p_j\), interpreted in the extended-real sense with \(0\log 0=0\) and \(u_j\log 0=-\infty\) when \(u_j>0\). Let \(\widehat p(u)=u/N\), set \(T_p(u)=2\{\ell_{\widehat p(u)}(u)-\ell_p(u)\}\), and define
\[
 \operatorname{pv}_N(p;u)
 =\sum_{z:\,\sum_j z_j=N,\ T_p(z)\geq T_p(u)}
   \operatorname{Mult}_N(z;p),
 \qquad
 \Gamma_N(u;\gamma)=\{p\in\Delta:\operatorname{pv}_N(p;u)\geq\gamma\}.
\]
For \(N=0\), the sole count vector is \(u=0\); define \(T_p(0)=0\), \(\operatorname{pv}_0(p;0)=1\), and \(\Gamma_0(0;\gamma)=\Delta\). For each arm, use the multinomial categories \(\{\varnothing\}\cup(\mathcal T\times\mathcal V)\), with probabilities \(p^a_{tv}=b_{tav}\) and \(p^a_{\varnothing}=1-\sum_{t,v}b_{tav}\). For each period of the external sample, use categories \(\mathcal V_0\) with probabilities \(q_t\). Let \(\mathcal G_{n,m}(\alpha)\) contain every \((b',q')\) accepted by all four confidence sets at threshold \(\alpha/4\), conditional on the realized arm sizes for the two armwise constructions. At fixed \(\delta\), define
\[
 \mathcal C_{n,m,\delta}^{\mathrm{exact}}
 =\operatorname{cl}_{[0,+\infty]}
   \bigcup_{\substack{(b',q')\in\mathcal G_{n,m}(\alpha):\\
                       \mathcal P_\delta(b',q',\omega)\neq\varnothing}}
   \Theta_{I,\delta}(b',q',\omega).
\]
The enumeration includes zero observed cells and zero candidate cells, and it uses neither studentization nor a selected endpoint optimizer.

%%% FIELD exact-whole-set-theorem-statement
Fix \(\alpha\in(0,1)\), a prespecified \(\delta\in[0,1]\), and the policy mixture \(\omega\). Use the convention \(\Gamma_0(0;\gamma)=\Delta\) in \(\text{def:exact-whole-set-projection}\). Under the stated sampling model, for every \(P\in\mathfrak M_\delta(b,q,\omega)\) and every joint law of the trial and external count vectors having the stipulated armwise conditional and periodwise marginal multinomial laws,
\[
 \Pr_{\mathrm{samp}}\!\left\{
   \Theta_{I,\delta}(b,q,\omega)
   \subseteq \mathcal C_{n,m,\delta}^{\mathrm{exact}}
 \right\}\geq 1-\alpha .
\]
This probability statement is unconditional over the realized treatment-arm sizes and external counts and requires no independence among the four count vectors. It includes zero observed cells, an empty realized arm, and every finite, \(+\infty\), and \([0,+\infty]\) branch of the extended ratio correspondence. For every realization, \(\mathcal C_{n,m,\delta}^{\mathrm{exact}}\) is the closure of the projection onto the ratio coordinate of a finite union of explicitly enumerable feasibility systems: each system consists of finitely many likelihood-ratio ordering inequalities, the candidate simplex constraints, \(x\in\mathcal P_\delta(b',q',\omega)\), and the three-branch relation generated by \(\rho_{11}(x,\omega)\) and \(\rho_{21}(x,\omega)\). Consequently, an empty candidate fiber contributes no projected point, and no endpoint optimizer is selected.

%%% FIELD exact-whole-set-theorem-proof
Put \(\gamma=\alpha/4\). We first prove the finite-sample property of one inverted multinomial test. For a dimension \(d\) and total \(N\), define the finite count space
\[
 \mathcal Z_{N,d}=\left\{z\in\mathbb N_0^d:\sum_{j=1}^d z_j=N\right\}.
\]
When \(N\geq1\), let \(Z\sim\operatorname{Mult}_N(p)\) and use the extended likelihood convention in \(\text{def:exact-whole-set-projection}\). The statistic \(T_p(Z)\) has finitely many distinct values, which we order as \(\tau_1>\cdots>\tau_L\). Write
\[
 c_\ell=\Pr_p\{T_p(Z)\geq\tau_\ell\},\qquad \ell=1,\ldots,L.
\]
Then \(c_1\leq\cdots\leq c_L=1\), and on \(\{T_p(Z)=\tau_\ell\}\) the definition gives \(\operatorname{pv}_N(p;Z)=c_\ell\). The set of indices for which \(c_\ell<\gamma\) is either empty or an initial segment \(\{1,\ldots,r\}\). Hence, in the nonempty case,
\[
 \Pr_p\{\operatorname{pv}_N(p;Z)<\gamma\}
 =\sum_{\ell=1}^r\Pr_p\{T_p(Z)=\tau_\ell\}
 =c_r<\gamma,
\]
and the probability is zero in the empty case. Thus
\[
 \Pr_p\{p\in\Gamma_N(Z;\gamma)\}
 =\Pr_p\{\operatorname{pv}_N(p;Z)\geq\gamma\}
 \geq1-\gamma. \tag{1}
\]
This argument allows coordinates of \(p\) or \(Z\) to be zero: a positive count paired with \(p_j=0\) gives extended statistic \(+\infty\), while \(0\log0=0\). If \(N=0\), the corrected definition gives \(\Gamma_0(0;\gamma)=\Delta\), so (1) holds with probability one.

For arm \(a\in\mathcal A\), let \(N_a\) be its realized size and let \(U^a\) be its count vector on \(\{\varnothing\}\cup(\mathcal T\times\mathcal V)\). Define
\[
 p^a_{tv}=b_{tav},\qquad
 p^a_{\varnothing}=1-\sum_{t\in\mathcal T}\sum_{v\in\mathcal V}b_{tav}.
\]
The stated sampling model says that, conditional on \(N_a=r\), \(U^a\sim\operatorname{Mult}_r(p^a)\). Applying (1), including its \(r=0\) case, and then the law of total probability yields
\[
 \Pr_{\mathrm{samp}}\{p^a\notin\Gamma_{N_a}(U^a;\gamma)\}
 =\sum_{r=0}^n\Pr(N_a=r)
   \Pr\{p^a\notin\Gamma_r(U^a;\gamma)\mid N_a=r\}
 \leq\gamma. \tag{2}
\]
For external period \(t\in\mathcal T\), let \(W^t\) be the count vector on \(\mathcal V_0\). Its marginal law is \(\operatorname{Mult}_m(q_t)\), so (1) gives
\[
 \Pr_{\mathrm{samp}}\{q_t\notin\Gamma_m(W^t;\gamma)\}\leq\gamma. \tag{3}
\]
Let \(E\) be the event that both arm vectors and both external-period vectors are accepted. Equations (2)--(3) and the union bound give
\[
 \Pr_{\mathrm{samp}}(E^c)
 \leq\sum_{a\in\mathcal A}\gamma+\sum_{t\in\mathcal T}\gamma
 =4\gamma=\alpha. \tag{4}
\]
No factorization of the joint sampling law was used in (4). Notice also why (4) is asserted unconditionally: the sampling model supplies the armwise multinomial statement conditional on \(N_a\), but it supplies only the needed marginal statements for the external vectors and imposes no condition making their coverage remain valid after conditioning jointly on the realized arm sizes.

We next verify that the true fiber is nonempty. Fix \(P\in\mathfrak M_\delta(b,q,\omega)\), and set \(x_s=P(S=s)\). Then \(x_s\geq0\) and \(\sum_sx_s=1\). By \(\text{ass:no-vaccine-effect-on-exposure}\), \(x_s=0\) whenever \(h^0\neq h^1\). For every \((t,a,v)\in\mathcal T\times\mathcal A\times\mathcal V\), \(\text{ass:history-indexed-first-infection-map}\), \(\text{ass:trial-randomization}\), \(\text{ass:arm-positivity}\), and \(\text{ass:outcome-consistency}\) imply
\[
 \begin{aligned}
 \sum_s\widetilde g_{tav}(s)x_s
 &=\Pr_P(J_t^a=v)\\
 &=\Pr_P(J_t^a=v\mid A=a)\\
 &=\Pr_P(J_t=v\mid A=a)
 =b_{tav}. \tag{5}
 \end{aligned}
\]
Here the second equality uses randomization and the conditioning event has positive probability by arm positivity. By \(\text{ass:external-margin-transport}\),
\[
 \sum_s\mathbf 1\{h_t^0=v\}x_s=q_{tv}
 \quad\text{for every }(t,v)\in\mathcal T\times\mathcal V_0. \tag{6}
\]
The equality in \(\text{ass:placebo-reference-stability}\), expanded under \(x_s=P(S=s)\), is
\[
 \sum_{s,v}\omega_v\{r_{10v}(s)-r^{[0]}_{20v}(s)\}x_s=0. \tag{7}
\]
Finally, for each \((a,u,v)\in\mathcal A\times\mathcal V\times\mathcal V\), \(\text{ass:bounded-noninfecting-exposure-priming}\) is exactly
\[
 \begin{aligned}
 &\sum_s\mathbf 1\{h_1^0=u,r_{1au}=0\}
          \mathbf 1\{r^{[u]}_{2av}\neq r^{[0]}_{2av}\}x_s\\
 &\hspace{35mm}\leq
 \delta\sum_s\mathbf 1\{h_1^0=u,r_{1au}=0\}x_s. \tag{8}
 \end{aligned}
\]
Equations (5)--(8), the support restriction, nonnegativity, and normalization are precisely all constraints in \(\text{def:reduced-history-polytope}\). Therefore
\[
 x\in\mathcal P_\delta(b,q,\omega),
 \qquad
 \mathcal P_\delta(b,q,\omega)\neq\varnothing. \tag{9}
\]

On \(E\), the true pair \((b,q)\) belongs to \(\mathcal G_{n,m}(\alpha)\). By (9), it is not removed by the nonempty-fiber condition. The definition of the projection therefore gives the deterministic inclusion
\[
 \Theta_{I,\delta}(b,q,\omega)
 \subseteq
 \bigcup_{\substack{(b',q')\in\mathcal G_{n,m}(\alpha):\\
                     \mathcal P_\delta(b',q',\omega)\neq\varnothing}}
 \Theta_{I,\delta}(b',q',\omega)
 \subseteq\mathcal C_{n,m,\delta}^{\mathrm{exact}}. \tag{10}
\]
Combining \(E\subseteq\{\Theta_{I,\delta}(b,q,\omega)\subseteq\mathcal C_{n,m,\delta}^{\mathrm{exact}}\}\) from (10) with (4) proves the claimed coverage. By \(\text{thm:extended-sharp-identified-set}\), the set on the left side of (10) is the sharp compatible-law ratio set, rather than merely an outer functional. In particular, its positive-denominator, \(+\infty\), and \([0,+\infty]\) cases are all covered simultaneously; no division by a possibly zero denominator occurred in the coverage argument.

It remains to give the promised finite representation. Index the four tests by
\[
 \mathcal I=\{(\mathrm{tr},0),(\mathrm{tr},1),(\mathrm{ex},1),(\mathrm{ex},2)\}.
\]
For test \(i\), let \(u_i\), \(N_i\), and \(d_i\) denote its observed count vector, total, and number of categories. Thus \(N_{(\mathrm{tr},a)}=N_a\), \(N_{(\mathrm{ex},t)}=m\), \(d_{(\mathrm{tr},a)}=1+2K\), and \(d_{(\mathrm{ex},t)}=K+1\). Let \(p_i(b',q')\) be the corresponding affine candidate vector: it is \((p^{a\prime}_{\varnothing},(b'_{tav})_{t,v})\) for a trial arm and \(q'_t\) for an external period.

For \(N_i\geq1\) and every subset \(B_i\subseteq\mathcal Z_{N_i,d_i}\), define
\[
 \begin{aligned}
 D_i(B_i)=\biggl\{p\in\Delta_{d_i}:{}&
 \bigl[T_p(z)\geq T_p(u_i)\iff z\in B_i\bigr]
       \text{ for every }z\in\mathcal Z_{N_i,d_i},\\
 &\sum_{z\in B_i}\operatorname{Mult}_{N_i}(z;p)\geq\gamma
 \biggr\}. \tag{11}
 \end{aligned}
\]
For \(N_i=0\), take the single label \(B_i=\mathcal Z_{0,d_i}=\{0\}\) and put \(D_i(B_i)=\Delta_{d_i}\). Since the count spaces are finite, there are finitely many labels \(B_i\). Directly from the definition of the tail probability,
\[
 \Gamma_{N_i}(u_i;\gamma)=\bigcup_{B_i\subseteq\mathcal Z_{N_i,d_i}}D_i(B_i) \tag{12}
\]
when \(N_i\geq1\), with the just-stated singleton interpretation when \(N_i=0\). Indeed, a candidate \(p\) selects the unique tail set \(B_i(p)=\{z:T_p(z)\geq T_p(u_i)\}\), and its tail sum is the second expression in (11).

Define the extended ratio relation \(\mathscr R\subseteq[0,+\infty]\times[0,1]^2\) by
\[
 \begin{aligned}
 (r,r_1,r_2)\in\mathscr R\quad\Longleftrightarrow\quad
 &\bigl[r_2>0\ \text{ and }\ r=r_1/r_2\bigr]\\
 &{}\text{ or }\ \bigl[r_2=0<r_1\ \text{ and }\ r=+\infty\bigr]\\
 &{}\text{ or }\ \bigl[r_1=r_2=0\ \text{ and }\ r\in[0,+\infty]\bigr]. \tag{13}
 \end{aligned}
\]
For a tuple \(\mathbf B=(B_i:i\in\mathcal I)\), let \(\mathfrak F_{\mathbf B}\) be the set of all \((r,b',q',x)\) satisfying
\[
 \begin{gathered}
 p_i(b',q')\in D_i(B_i)\quad(i\in\mathcal I),
 \qquad x\in\mathcal P_\delta(b',q',\omega),\\
 \bigl(r,\rho_{11}(x,\omega),\rho_{21}(x,\omega)\bigr)\in\mathscr R. \tag{14}
 \end{gathered}
\]
The simplex restrictions on \((b',q')\) are included in the first line of (14). Substituting (12) and (13) into \(\text{def:exact-whole-set-projection}\) gives the exact equality
\[
 \mathcal C_{n,m,\delta}^{\mathrm{exact}}
 =\operatorname{cl}_{[0,+\infty]}
   \left(
     \bigcup_{\mathbf B\in\prod_{i\in\mathcal I}\mathfrak B_i}
     \operatorname{proj}_r\mathfrak F_{\mathbf B}
   \right), \tag{15}
\]
where \(\mathfrak B_i=2^{\mathcal Z_{N_i,d_i}}\) for \(N_i\geq1\) and \(\mathfrak B_i=\{\mathcal Z_{0,d_i}\}\) for \(N_i=0\). Formula (15) is finite: every \(\mathfrak B_i\) and every count space is finite; (11) contains finitely many likelihood-order comparisons and one finite multinomial sum; membership in \(\mathcal P_\delta\) is the displayed finite linear system; and (13) has three branches. If \(\mathcal P_\delta(b',q',\omega)=\varnothing\), then no \(x\) satisfies (14), so that candidate contributes nothing exactly as required. Conversely, every accepted candidate and every \(x\) in its nonempty fiber selects one tuple \(\mathbf B\) and one branch of (13), so no point is lost. Thus (15) is an exact feasibility projection and never selects, or conditions on, an endpoint optimizer.
